\section{Challenges of IoT Edge Caching}

Edge caching faces challenges due to changing demand, user mobility, limited storage and computing resources, and evolving relationships between nodes. Demand varies across locations and over time, and user movement creates correlated and non-stationary request patterns that do not match classical caching assumptions. Limited edge resources and increasing network density also make coordination and scalability more difficult. These factors make caching at the IoT edge more complex than in traditional systems and require strategies that consider mobility and network structure.


IoT workloads exhibit strong spatial and temporal variability, where data popularity differs across locations and changes over time due to contextual or event driven factors~\cite{Satyanarayanan2017Edge,Chiang2016Fog}. User mobility further creates correlation across edge nodes, as content requested at one access point may soon be requested at neighboring nodes when users move \cite{Kotz2004Dartmouth}. At the same time, the arrival and departure of users can cause sudden changes in local popularity distributions. These effects challenge classical caching models, which typically assume stable, independent, and stationary request streams \cite{Podlipnig2003}


Edge nodes such as IoT gateways or access points operate under constrained storage budgets. Unlike centralized cloud servers or CDNs, edge devices must carefully allocate limited cache capacity among competing content. This creates a trade-off between replicating content across nearby nodes to increase availability and diversifying cached content to reduce redundancy. In addition to storage constraints, IoT edge devices often have limited computational power, energy resources, and communication bandwidth which limit the feasibility of fully centralized optimization and motivate scalable, distributed approaches. Furthermore, although the physical infrastructure of edge networks may remain fixed, the logical interaction structure among nodes evolves over time due to user mobility and varying connectivity condition. So, caching strategies must therefore adapt not only to changing demand but also to evolving inter-node relationships. Also as future 5G/6G and IoT deployments become increasingly dense, the number of potential cooperative relationships grows, making coordination and optimization more complex and less practical. Any effective caching strategy must therefore scale efficiently with network size while maintaining decentralized operation.

Collectively, these challenges motivate IoT edge caching strategies that consider mobility-induced relationships and adapt to dynamic demand under limited resources. Effective solutions must consider how nodes interact over time and coordinate caching decisions in a scalable manner.

